% Generated draft: requires mathematical and source review.
\documentclass[11pt]{article}
\usepackage[margin=25mm]{geometry}
\usepackage{fontspec}
\setmainfont{Latin Modern Roman}
\usepackage{amsmath,amssymb,mathtools,unicode-math}
\setmathfont{Latin Modern Math}
\usepackage{enumitem,needspace,xurl,hyperref}
\setlength{\parindent}{0pt}
\setlength{\parskip}{5pt}
\setlength{\emergencystretch}{3em}
\newcommand{\N}{\mathbb{N}}
\newcommand{\Z}{\mathbb{Z}}
\newcommand{\Q}{\mathbb{Q}}
\newcommand{\R}{\mathbb{R}}
\newcommand{\C}{\mathbb{C}}
\newcommand{\F}{\mathbb{F}}
\newcommand{\A}{\mathbb{A}}
\newcommand{\PP}{\mathbb{P}}
\newcommand{\E}{\mathbb{E}}
\newcommand{\eps}{\varepsilon}
\newcommand{\dd}{\,\mathrm{d}}
\begin{document}
\section*{UnsolvedMath ID: 7000002}
\textbf{Source problem code:} catalogue entry
\subsection{Definitions and mathematical statement}
In a standard compact-surface interpretation, let $f:M^2\to\R^3$ be a smooth immersion of a closed connected surface. One common tightness convention requires every generic height function $h_v(x)=\langle f(x),v\rangle$, $v\in S^2$, to be a perfect Morse function over $\Z/2$: the number of critical points of each index equals the corresponding Betti number. The source does not give its exact tightness definition or state whether surfaces are embedded, immersed, or necessarily compact.

Under a deformation-rigidity interpretation, the question would ask whether every smooth isometric family $f_t$ with $f_0=f$ satisfies
\[
f_t^*g_{\mathrm{Euc}}=f^*g_{\mathrm{Euc}}quad\Longrightarrow\quad
f_t=Q_t f+a_t,
\qquad Q_t\in O(3),\ a_t\in\R^3.
\]
Here the conclusion is that every deformation is induced by ambient rigid motions. Global rigidity, asserting congruence of every isometric realization, and infinitesimal rigidity, concerning first-order variations, are different statements. The source's single word rigid does not distinguish them or say whether compared realizations must themselves remain tight.

The literal question is whether all smooth tight surfaces in $\R^3$ are rigid in the intended sense. No convexity, orientability, genus or embedding assumption is added. The supplied summary reports classical convex rigidity and describes the nonconvex smooth question as unresolved. Its vague pinched/nonrigid examples and projective-plane language do not identify an admissible smooth embedded counterexample; those category distinctions need checking.
\subsection{Short English statement}
Are smooth tight surfaces in three-space rigid under isometric comparison or deformation? The source must specify tightness, immersion versus embedding, compactness and the exact notion of rigidity.
\subsection{Sources}
\begin{itemize}
\item Local UnsolvedMath registry, permanent ID 7000002, exact source code catalogue entry, ``Geometry of Curves and Surfaces---Problem 1.2''; proposed by Mohammad Ghomi, 2019. Registry status open, classification OPEN-TRIAGE.
\item Mohammad Ghomi, \emph{Open Problems in Geometry of Curves and Surfaces} (2019), Problem 1.2, PDF page 6, \url{https://people.math.gatech.edu/~ghomi/Papers/op.pdf}. Supplied normalized pdf-text extraction accessed July 29, 2026. No PDF was accessed here.
\item The supplied summary names Cohn--Vossen for classical convex rigidity, discusses nonconvex tight surfaces, and says it located no definitive general solution during 2019--2026. It mentions pinched examples and a tight projective plane without precise references or category hypotheses. In particular an immersed projective plane must not be silently treated as an embedded surface in $\R^3$.
\item Provenance flags current status and statement-reuse terms NEEDS\_REVIEW; the source is publicly author-hosted.
\end{itemize}
\end{document}
